\chapter{Representing messages: the numerical layer}
\label{ch:numerics}

Chapter~\ref{ch:bp} gave an exact algorithm on functions. A computer cannot store a function,
so an implementation must choose a representation --- and that choice is the \emph{only}
approximation in the whole scheme. This chapter is about it, because conflating this
approximation with the exactness of the previous chapter is the easiest way to overstate what
the method does.

\section{Three statements that must be kept apart}
\label{sec:three-statements}

\begin{enumerate}
\item \textbf{Sum-product on the posterior chain is exact}, as a theorem about tree-structured
factorisations (Proposition~\ref{prop:bp-exact}). This is a mathematical statement about
functional messages and involves no computer.
\item \textbf{The discretised recursion is exact for the discretised model.} Restricting
messages to a grid defines a different, finite-state model; on that model the recursion is
exact up to floating-point arithmetic.
\item \textbf{The discretised model approximates the continuous one}, with an error that must
be \emph{measured}.
\end{enumerate}

\begin{pitfall}
Writing simply ``our method is exact'' collapses (1) and (3) and claims something false.
Throughout this project we write \emph{exact tree inference} for (1) and \emph{exact for the
discretised model} for (2), and never abbreviate. The distinction is not pedantry: (1) holds
for any innovation law whatsoever, while the error in (3) depends strongly on the smoothness
of that law --- as \S\ref{sec:quadrature} shows.
\end{pitfall}

\section{The grid}

Messages are represented by their values on $\ngrid$ equally spaced points on $[-A,A]$, and
integrals become composite trapezoidal sums,
\begin{equation}
\int\!\mathrm{d}a\, f(a) \;\approx\; \sum_{k=1}^{\ngrid} w_k f(u_k),
\qquad
w_k = \Delta u \ \text{(interior)},\quad \tfrac{1}{2}\Delta u \ \text{(endpoints)} .
\end{equation}
Each message update \eqref{eq:fwd} becomes a matrix--vector product with the matrix
$\kernel(u_l\mid u_k)$, which is formed once per model and reused at every site and every
sequence.

\begin{codebox}
\coderef{src/bp\_grid.py}{make\_grid} builds the points and trapezoidal weights.
The transition matrix comes from each prior's \texttt{log\_transition\_matrix(grid)}
(\codefile{src/priors.py}) or each estimated kernel's (\codefile{src/kernels.py}) --- the same
interface, which is why the true and learned models are interchangeable in every experiment.
\end{codebox}

\section{Two error sources}

\subsection{Truncation}

Probability mass outside $[-A,A]$ is discarded. The diagnostic is the mass in the boundary
cells: if it is not negligible, $A$ is too small and refining $\ngrid$ will not help, because
the error is in the domain rather than the resolution.

\begin{pitfall}
This failure is invisible unless one looks for it. An earlier package in this project ran at
$A=8$ with a boundary mass of $7.4\times10^{-7}$, above its own tolerance, and the symptom was
not an error message but a slow drift in the posterior means at large $t$. Widening the domain
at preserved spacing reduced it to $5.3\times10^{-10}$.
\end{pitfall}

\subsection{Quadrature}
\label{sec:quadrature}

The trapezoidal rule has error $O(\Delta u^2)$ for integrands with a bounded second
derivative. Whether that bound applies depends on the innovation.

\begin{itemize}
\item \textbf{Gaussian innovation.} The integrand is smooth and convergence is fast.
\item \textbf{Laplace innovation.} $\innov(e)\propto e^{-|e|/b}$ has a discontinuous derivative
at $e=0$. The second derivative is not bounded there, the standard estimate does not apply, and
convergence is second order at best.
\end{itemize}

\begin{pitfall}
It is tempting --- and this project did it in an earlier draft --- to describe the
discretisation as ``spectrally accurate'' on the strength of the Gaussian measurements. That
generalisation is not available: spectral accuracy requires smoothness the Laplace kernel does
not have. The accompanying note now claims second-order behaviour for that family and nothing
more.
\end{pitfall}

\subsection{A resolution condition}

At small $t$ the likelihood $\like(x_i\mid a_i)$ is a narrow Gaussian of width
$\sqrt{\Dt}/e^{-t}$. If that width is comparable to the grid spacing, the integrand is
under-resolved and the quadrature is meaningless however fine the rest looks. The working
condition used here is $\sqrt{2t} \gtrsim 3\Delta u$.

\begin{intuition}
Everything gets harder as $t\to0$, from two directions at once: the likelihood narrows, and the
score \eqref{eq:score-final} carries a $1/\Dt$ that diverges. This is why reverse integration
must stop at some $t_{\min}>0$ and why the small-$t$ end of every table deserves the most
suspicion.
\end{intuition}

\section{Stability}

Messages are products of many small numbers and underflow quickly. Two standard devices:
work in the log domain where possible with per-row maximum subtraction, and renormalise each
message to unit quadrature mass after every step, accumulating the discarded constants so the
evidence is recovered exactly.

\begin{codebox}
Both appear in \coderef{src/bp\_grid.py}{grid\_bp}: \texttt{log\_ell -= log\_ell.max(...)}
before exponentiating, and an explicit mass check that raises \texttt{FloatingPointError} if a
message loses all its mass rather than silently returning garbage.
\end{codebox}

\section{The stopping-time problem}
\label{sec:tmin}

Because integration stops at $t_{\min}>0$, what a sampler produces is a draw from
$P_{t_{\min}}$, not from $\Pz$. This sounds like a technicality and is not.

Independent Gaussian noise of variance $v$ added to a quantity of variance $\ivar$ multiplies
its excess kurtosis by $(\ivar/(\ivar+v))^2$. At $\corr=0.85$ and $t_{\min}=0.02$ this predicts
a generated excess kurtosis of $1.910$ where the clean chain has $3.0$.

\begin{pitfall}
This project measured exactly $1.91$ for the reference sampler and initially read it as a
failure of the integrator to converge. It was not: the sampler was reproducing
$P_{t_{\min}}$ correctly and being compared against the wrong target.

The obvious repair --- apply a final denoising step and report $\pmean{a}$ instead of $x$ ---
\emph{inverts} the bias rather than removing it. The posterior mean of a heavy-tailed prior is
a shrinkage operator: it pulls small values towards zero harder than large ones, so the law of
the posterior mean is \emph{more} peaked than $\Pz$. Measured overshoot: $3.73$, $4.34$, $4.74$
at $t_{\min}=0.02,0.05,0.10$ against a true $3.0$ --- and identical at $\ngrid=401$ and
$\ngrid=801$, which rules out grid resolution as the cause.

The correct treatment is to compare like with like: draw from $\Pz$, noise it forward to the
same $t_{\min}$, and compare. Both sides are then samples of the same law, no estimator sits
between the sampler and the metric, and the target moments follow in closed form.
\end{pitfall}

\begin{codebox}
\coderef{experiments/exp\_16\_sampling\_validation.py}{reference\_at} builds the forward-noised
reference; \texttt{target\_innovation\_moments} gives the closed-form target. The docstrings
there record both failed designs, since the reasoning is easier to follow than to rediscover.
\end{codebox}

\section{Running on a GPU}

The recursion is a sequence of dense matrix products, which suits a GPU. Rather than write a
second implementation, \codefile{src/backend.py} parameterises the existing one by its array
module, so CPU and GPU execute the same code.

\begin{intuition}
Two implementations of an exact algorithm is a bad trade. They agree on the easy cases and
drift where nobody is looking, and the exactness claims of \S\ref{sec:bp-validation} would then
hold for one device and be assumed for the other.
\end{intuition}

\begin{codebox}
\codefile{tests/test\_backend\_parity.py} gates every GPU number: device agreement, the
closed-form Gaussian check re-run on device, float64 preservation, and invariance to batch
width. Measured on an NVIDIA H200: posterior means agree to $1.8\times10^{-16}$--$8.0\times10^{-16}$,
variances to $9.3\times10^{-15}$, and the GPU reproduces the closed-form Gaussian score to
$4.2\times10^{-16}$. \coderef{hpc/bocconi\_gpu.sbatch}{} runs that suite before any sweep and
exits if it fails.
\end{codebox}

\begin{pitfall}
The availability probe originally tested that a GPU array could be \emph{allocated}. Allocation
succeeds without cuBLAS loaded, so on a misconfigured node the check passed and every matrix
product then failed inside the recursion. It now probes an actual matrix product. A capability
check must exercise the operation that will be used, not a cheaper proxy.
\end{pitfall}
